\documentclass[10pt]{article}
\usepackage[margin=0.8in]{geometry}
\usepackage{amsmath,amssymb}
\usepackage[T1]{fontenc}
\usepackage[utf8]{inputenc}
\usepackage{lmodern}
\usepackage{microtype}
\usepackage{enumitem}
\usepackage{booktabs}
\usepackage{array}
\usepackage{url}
\usepackage[hidelinks]{hyperref}
\def\UrlBreaks{\do\/\do\_\do-}
\setlist{nosep,leftmargin=1.3em}
\setlength{\parskip}{3pt}\setlength{\parindent}{0pt}
\sloppy
\newcommand{\fpath}[1]{\nolinkurl{#1}}
\newcolumntype{L}[1]{>{\raggedright\arraybackslash}p{#1}}
\begin{document}

\begin{center}
{\Large\bfseries Measuring the MOND Acceleration Scale at High Redshift: What Current Public Data Can and Cannot Say (the Baryon-Calibration Wall)}\\[6pt]
{\small The authors --- draft for review, 2026-10-01, version 1.2 (wording corrections after an independent referee, CFG241; no result changes; one arithmetic slip corrected, F-14) --- AI-assisted research programme; not peer reviewed}
\end{center}

% Conventions of this draft. Every number carried by an audit row of PAPER38_audit.py is checked against a committed file (git HEAD)
% and against the tex block that carries its tag; numbers quoted only in running text are not checked by the script (see the Reproducibility section).
% A line "% AUDIT: Bnn" after a paragraph, item or table row tags the audited numbers quoted in the text just above it.
% v1.2: tags renumbered B01-B43 in text order (v1.1 had no B34 and placed B43 after B14; v1.1 B43 -> B15, B15-B33 -> B16-B34, B42 -> B43; B42 is now the Reproducibility block).
% No new analysis is done here: this note compiles committed lanes; where a lane README and its script output differ, the output wins.

\begin{abstract}
\noindent The programme ties the MOND acceleration scale to the vacuum, $a_0=\kappa c\sqrt{G\rho_\Lambda}$, with $\kappa=\tfrac12$ fitted, not derived. Its distinctive high-redshift prediction is a flat $a_0(z)$ (distinctive against $\Lambda$CDM and the $H(z)$ rival; standard MOND shares it); the rival is $a_0\propto H(z)$, about $\times2.2$ at $z\approx1.4$ and about $\times8$ at $z\approx5$; $\Lambda$CDM has no $a_0$ and enters only through an effective-$a_0$ proxy. We collect the programme's committed high-redshift $a_0$ lanes and state one result plainly: in every lane the limit that survives an unlimited sample is the absolute calibration of the baryon (gas and stellar) masses; at present sizes the KiDS split, the class-S sources of the compilation and the seven class-M galaxies are also limited by statistical power or by the Newtonian regime (CFG255, CFG227, CFG229). (i)~Where the discs sit (per-point median $g_{\rm bar}/a_0=1.1$ to 4.4) a 0.05\,dex error in the mass discrepancy is a factor 1.5 to 1.8 in $a_0$, and the implied $a_0$ moves by $-2.5$ to $-4.8$ dex per dex of baryon mass (a post hoc finite-difference estimate at a $\pm0.03$\,dex step; the step changes the range, $-2.5$ to $-6.4$ at $\pm0.01$\,dex); in the deep regime an amplitude test measures only the product of the calibration and $a_0$. (ii)~No sample separates the laws: RC100 is gas-route-limited, CRISTAL is route-dependent and calibration-limited, and the $z=2$--5 compilation, the ALMA [CII] rotators and the seven class-M galaxies are NOT POSSIBLE at pre-flight. (iii)~The gas calibration at $z\approx2.2$ is a prescription bracket (${\sim}0.2$--0.7\,dex); tracer agreement does not fix the absolute scale, and no public dynamics-independent anchor reaches 0.10\,dex (the best route, an absorption-line $\alpha_{\rm [CI]}(Z)$ relation, does not reproduce its stated 0.2\,dex scatter from our transcription of its Table~1: 0.64\,dex about the quoted relation; how the authors define that scatter, and any erratum, were not checked). (iv)~A within-survey lens-redshift split, which would cancel a common-mode calibration, is NOT POSSIBLE in KiDS-1000 (power 0.14 against the 9 required) and NON-DISCRIMINATING by its own MUTATE control; even with unlimited statistics a stellar-mass drift $\delta$ would mimic an $a_0$ drift by $\delta/2$. A decisive $a_0(z)$ test needs baryon errors and band $\le0.10$\,dex. With the public data on hand the question is open; nothing here is a detection either way. The near-term test of the programme is a separate, $z=0$ one, which its current candidate law (B) can only survive or fail: Gaia DR4 wide binaries (2 December 2026).
\end{abstract}
% AUDIT: B01

\section{The prediction, the rival, and the size of the signal}\label{sec:pred}

\textbf{The law and its footings.} The programme's galaxy law is a MOND kernel $\nu$ (the lanes use the monotone interpolation $\nu_{\rm mono}$; the CFG240 theorems of Section~\ref{sec:lever} use the exact P2 kernel) with $a_0=\kappa c\sqrt{G\rho_\Lambda}$, $\kappa=\tfrac12$ fitted (not derived) and the cold component's mass still required; no dark-matter particle is added. Two footings of the local value are carried throughout: canonical $a_0=9.3603\times10^{-11}$ and alternative $1.1312\times10^{-10}$\,m\,s$^{-2}$. The canonical footing decides unless stated; the implied ratio $s^*$ of Section~\ref{sec:lever} is footing-relative (the alternative footing scales it by $1/1.209$, the absolute $a_0$ being unchanged).
% AUDIT: B02

\textbf{Three readings of $a_0(z)$.} \emph{Flat} (the programme's distinctive law, distinctive against $\Lambda$CDM and the $H(z)$ rival; standard MOND shares a constant $a_0$): $a_0(z)=a_0(0)$. \emph{The rival}: $a_0(z)=a_0(0)\,E(z)$, with $E=H/H_0=\sqrt{\Omega_m(1+z)^3+1-\Omega_m}$ (the reading of the coincidence $a_0\sim cH$ as an evolving scale; the repository's lensing note ties this shape to Verlinde's emergent-gravity proposal). The lanes use $\Omega_m=0.30$ to 0.315, across which $E(1.4)=2.20$ to 2.25 and $E(5)=8.09$ to 8.29 (CFG241); the rival numbers below carry that few-per-cent dependence. \emph{The $\Lambda$CDM effective-$a_0$ proxy}: $\Lambda$CDM has no $a_0$; the proxy is a halo-structure scale of the record's making used as if it were $a_0(z)$ (Dutton \& Macci\`o concentrations [DM14]), $F=1.23$, 1.77 and 2.16 at $z=1$, 2 and 2.5 ($+0.33$\,dex at 2.5), 4.52 at $z=4.5$ and 6.20 at $z=5.5$. It is extrapolated above $z=5$, and a fair $\Lambda$CDM test uses simulated galaxies (CFG222).
% AUDIT: B03

\textbf{How large the rival's signal is, by regime.} The rival's $a_0$ grows with $z$, but what a test sees depends on where the galaxies sit in $g_{\rm bar}/a_0$: in the deep regime a change of $a_0$ shows at half its size in the amplitude, and near the Newtonian regime it hardly shows at all.
\begin{center}\small
\setlength{\tabcolsep}{4pt}
\begin{tabular}{L{0.23\linewidth}L{0.21\linewidth}L{0.36\linewidth}L{0.12\linewidth}}
\toprule
regime & rival / flat in $a_0$ & rival-vs-flat signal in the observable & lane \\
\midrule
KiDS lenses, $z\approx0.20\to0.40$, deep regime & small over the lever & amplitude change $+0.018$\,dex (flat $+0.001$) & CFG255 \\
% AUDIT: B04
NOEMA3D, $z\approx1.4$ & about $\times2.2$ & (both laws consistent there; Section~\ref{sec:samples}) & CFG213 \\
% AUDIT: B05
RC100, $z=0.61$--2.52 & --- & slope of $\delta_{\rm flat}$ on $z$: 0.000 if flat is true, $+0.073$ if the rival is & CFG216 \\
% AUDIT: B06
$z\approx2$ to 2.3 & 3.4 at $z=2.3$; a factor-3 law difference (0.48\,dex) at $z\approx2$ & implied-$a_0$ ratio $s^*$ & CFG223 \\
% AUDIT: B07
CRISTAL and ALPINE, $z\approx5$ & 8.69 and 8.84 (the two CRISTAL points) & median $\delta$ separation 0.288 (CRISTAL) and 0.287 (ALPINE6); 0.94\,dex in $a_0$ & CFG218, CFG228, CFG224 \\
% AUDIT: B08
class-M galaxies, $z\ge2$, $g_{\rm bar}/a_0=4.8$ to 157 & as above & 0.005 to 0.09\,dex per galaxy & CFG229 \\
% AUDIT: B09
\bottomrule
\end{tabular}
\end{center}
The largest observable signal, 0.29\,dex at $z\approx5$ (CFG218), is smaller than the joint outer-band baryon systematic of 0.498\,dex that the ALMA lane carries on the same statistic (CFG228). That comparison is the subject of this note.
% AUDIT: B10

\section{The lever: why a small baryon error is a large $a_0$ error}\label{sec:lever}

\textbf{The implied $a_0$.} CFG223 defines the implied ratio $s^*$ as the multiplicative $a_0$ scale (relative to the footing's local value) that sets the median over a set of galaxies of $\delta_i(s)=\log_{10}[D_i/\nu(g_{{\rm bar},i}/(a_0s))]$ to zero, with $D=g_{\rm obs}/g_{\rm bar}$. Flat predicts $s^*=1$ only if the baryon calibration that set the local $a_0$ holds at every $z$. The statistical bars are small and the baryon-mass bands span the plot: RC100's 95\% intervals are a factor 2.2 to 3.8 wide, while a $\pm0.15$\,dex shift of all baryon masses moves $s^*$ from 0.45 to 4.04 in the lowest-$z$ quartile and removes the solution altogether (median $D$ below 1) in quartiles 3 and 4; $\pm0.30$\,dex removes it at every point.
% AUDIT: B11

\textbf{The lever.} The eight points have median $g_{\rm bar}/a_0=1.1$ to 4.4 (individual discs 0.19 to 18.2), where the kernel varies slowly (local slope 0.13 to 0.28 against 0.5 in the deep regime). There a 0.05\,dex error in $D$ is a factor 1.5 to 1.8 in $a_0$ (a 0.05\,dex step; stable), and $d\log_{10}s^*/d(\text{baryon dex})=-2.5$ to $-4.8$ (post hoc diagnostic; the committed output prints $-2.48$ to $-4.83$ over the eight points, from a finite difference of the median root at $\pm0.03$\,dex). This baryon lever depends on the step: re-evaluated at $\pm0.01$\,dex the eight values run $-2.46$ to $-6.42$ (CFG241), and two points with nearly the same local slope give $-4.37$ and $-2.65$; the range is an order-of-magnitude guide, not a measured coefficient. To put a factor-3 law difference (0.48\,dex, $H(z)$ against flat at $z\approx2$) at 2$\sigma$, the shared baryon calibration would have to be known to about 0.05 to 0.10\,dex (an order-of-magnitude statement from that table).
% AUDIT: B12

\textbf{The same lever in the other lanes.} In the ALMA lane (CFG228) the lever is $\lambda=-1.8$ at the measured pooled root and $-3.3$ at the nominal $s^*=1$: a 0.1\,dex baryon error moves $\log_{10}s^*$ by 0.18 to 0.33\,dex, against 0.66\,dex between the flat and proxy expectations and 0.87\,dex between flat and $H(z)$. In the class-M lane (CFG229) the one galaxy at $g_{\rm bar}/a_0=4.8$ (ALESS 122.1) has $\lambda=-2.0$; for the discs at $g_{\rm bar}/a_0>10$ the pre-flight lever is $-56$ to $-70$, so a 0.05\,dex baryon error would move $\log_{10}a_0$ by 2.8 to 3.5 (a factor 600 to 3,000).
% AUDIT: B13

\textbf{The deep regime.} Deep in the MOND regime $g_{\rm obs}\simeq\sqrt{g_{\rm bar}a_0}$, so an amplitude test there measures the product of the baryon calibration factor and $a_0$: in the KiDS lensing lane (CFG255) a stellar-mass calibration drift $\delta$ between two redshift bins moves the amplitude by $\delta/2$, as an $a_0$ drift does (exactly in the deep limit; to within 6.1\% in the KiDS K1 bins for the P2 kernel, CFG241). Outside the deep limit the kernel's shape does in principle separate the calibration from $a_0$, so $a_0$ is not unidentifiable in general; but the separation rests on the curvature of $\nu$, and the estimate becomes ill-conditioned, its error growing without bound as the sample's largest $g_{\rm bar}/a_0$ falls toward the deep regime. The formal treatment is CFG240 (below).
% AUDIT: B14

\textbf{The conditioning, as theorems (CFG240).} The lane \fpath{campaign_fresh_gravity/CFG240_calibration_wall/} formalises this in the ChainCert Lean library (50 theorems; the library goes from 281 to 331, standard axioms only). Lean certifies implications from stated premises, not any empirical fact. With one calibration factor $f$ on $g_{\rm bar}$:
\begin{itemize}
\item \emph{Deep limit.} $g_{\rm obs}$ depends on $(f,a_0)$ only through $f a_0$.
\item \emph{Identifiability.} For the exact P2 kernel, two distinct $g_{\rm bar}$ values determine $(f,a_0)$, and one value does not.
\item \emph{Conditioning.} The Jacobian determinant tends to zero when both points are deep (and, by the same bound, when both are Newtonian).
\item \emph{A floor.} With $f$ free, any sample whose local slopes $b_i$ lie in $[0,\tfrac12]$ (proved in Lean for P2; checked numerically, not in Lean, for $\nu_{\rm mono}$) and whose errors are independent Gaussians of size $\sigma$, with an invertible Fisher matrix, obeys $\sigma(\log a_0)\ge 3\sigma/\sqrt N$: at a per-point error $\sigma=0.1$\,dex, a 0.1\,dex $a_0$ needs $N>9$; at 0.2\,dex it needs $N>36$.
\end{itemize}
Its Fisher companion (log-uniform design, P2, $N=20$, $\sigma=0.1$\,dex, no prior on $f$) finds:
\begin{itemize}
\item A sample with $y_{\min}=0.01$ must reach $y=g_{\rm bar}/a_0\gtrsim8$ (the break-even on a 0.02\,dex grid of $y_{\max}$), and one with $y_{\min}=10^{-3}$ must reach $y\gtrsim5.5$.
\item With no deep points ($y_{\min}=0.1$) it never reaches 0.1\,dex: the best is 0.117, reached only by designs extending to $y$ of about 48.
\item Even at break-even the two estimates stay correlated ($\rho\approx-0.8$): the degeneracy is broken in $\sigma$, not removed.
\end{itemize}
The high-$z$ points of Section~\ref{sec:lever} have median $y=1.1$ to 4.4 with no deep points; a log-uniform design over that range is the regime where this table gives no 0.1\,dex leverage on $a_0$ without an external calibration of $f$. The individual discs spread wider: for per-galaxy errors of 0.1\,dex the formal leverage of the real per-galaxy $y$ would be 0.15 to 0.22\,dex per RC100 quartile and 0.089\,dex for the 100 RC100 galaxies pooled (CFG241, Fisher with $f$ free). The wall therefore rests on the per-galaxy error and on the common-mode $f$, not on the $y$ range alone.
% AUDIT: B15

\textbf{What the points show (descriptive).} With the caveats above (author decompositions, gas-route-limited, not a detection), flat $a_0$ (ratio 1) lies inside the 95\% statistical interval of 7 of the 8 CFG223 points; the exception, CRISTAL at $R_e$ on the fit route, sits 0.065\,dex below its interval, so the lane's frozen caption rule does not release the sentence ``consistent with a constant $a_0$''. Inside the statistical / $\pm0.15$ / $\pm0.30$\,dex band-widened intervals the counts are flat 7/8/8, proxy 5/7/7 and $H(z)$ 4/5/7. The independent-route CRISTAL points carry no information (one of them has no root in 10\% of resamples).
% AUDIT: B16

\section{What each sample gives}\label{sec:samples}

Each row is one committed lane; every lane froze its criteria before the numbers it scores. ``Fit route'' = the authors' dynamically fitted or prior-anchored baryon masses; ``independent'' = SED stellar mass plus a tracer gas mass. Class S/L/M = single tracer at local prescriptions / dust or [CII] / per-galaxy optimised multi-tracer gas (the class-M band is a band on tracer agreement, not on the absolute gas scale; CFG227 correction).
\begin{center}\small
\setlength{\tabcolsep}{3pt}
\begin{tabular}{L{0.09\linewidth}L{0.12\linewidth}L{0.08\linewidth}L{0.07\linewidth}L{0.14\linewidth}L{0.42\linewidth}}
\toprule
lane & sample & $z$ & $N$ & baryon route & verdict and result in one line \\
\midrule
CFG216, CFG217 & RC100 & 0.61--2.52 & 100 & fit, prior-anchored (gas scaling grows with $z$) & W-flat as frozen, but GAS-ROUTE-LIMITED: a differential baryon-mass systematic of about 0.15--0.25\,dex over $z$ 0.6--2.5 reverses it; it measures the differential baryon calibration, not $a_0(z)$, and its frozen significance is not to be quoted \\
% AUDIT: B17
CFG213 & NOEMA3D & ${\approx}1.4$ & 10 & independent (measured CO) & both laws CONSISTENT in every cell, on both mass routes; no separation \\
% AUDIT: B18
CFG214 & NOEMA3D outer radii & 1.1--1.6 & 10 & --- & GATE-LIMITED, NON-DIAGNOSTIC AS FROZEN: 1 of 10 passes the 5\% gate (beam-smeared model curves) \\
% AUDIT: B19
CFG213, CFG220 & ALMA-CRISTAL & 4.4--5.7 & 12 / 6 & fit (1-dex prior on $M_{\rm bary}$) / independent (dust) & fit route: rival DISFAVOURED-under, ROUTE-DEPENDENT; independent route at $R_{\rm out}$: BOTH-CONSISTENT, CALIBRATION-LIMITED (the class flips at a $+0.07$\,dex gas offset), not LOO-robust \\
% AUDIT: B20
CFG215 & four samples & 0.3--5.7 & --- & mixed & across-sample slope NON-DIAGNOSTIC by construction (route mixing; MUSE-DARK route bias $+0.512$\,dex) \\
% AUDIT: B21
CFG227 & $z=2$--5 RAR compilation & 2.0--4.4; 4.4--5.7 & 9 + 6 & class S / L & NOT separable in any scored group; 7 of the 9 class-S sources have $g_{\rm obs}<g_{\rm bar}$ \\
% AUDIT: B22
CFG228 & six ALPINE [CII] rotators + SPT0418-47 & 4.2--5.5 & 7 & class S/L (outer band $\pm0.671$\,dex) & blind pre-flight NOT POSSIBLE; pressure term alone moves the median $\delta_{\rm FLAT}$ from $-0.467$ to $+0.193$; implied $a_0$ UNINFORMATIVE \\
% AUDIT: B23
CFG229 & class-M gold & 2.024--2.935 & 7 & class M & pre-flight NOT POSSIBLE (Newtonian regime: signal 0.028 against SD 0.108); 6 of 7 have $D<1$ and no root; one galaxy (ALESS 122.1) gives $s^*=8.82$ \\
% AUDIT: B24
CFG255 & KiDS-1000 lenses & 0.20 / 0.40 & 181,477 & photometric $M_*$ & NOT POSSIBLE (power 0.14 against 9) and NON-DISCRIMINATING by its MUTATE \\
% AUDIT: B25
\bottomrule
\end{tabular}
\end{center}

\textbf{Reading the table.} \emph{RC100} (Nestor Shachar et al.\ [NS23]) is the tightest sample statistically, but its baryon masses are prior-anchored with a gas scaling that grows with $z$. On the attack lane's own axis the flat law's slope is exactly zero at a differential baryon-mass change of $-0.075$\,dex between $z\approx0.6$ and 2.5 and the rival's at $-0.25$\,dex; the two sit about 0.18\,dex apart, so a 3$\sigma$ separation needs that differential scale known to about $\pm0.06$\,dex on top of the statistics. Holding the gas fraction fixed with $z$ ($-0.251$ on that axis) or a ULIRG-like $\alpha_{\rm CO}$ ($-0.211$) gives the rival back; both sit at or beyond the lane's frozen $\pm0.2$\,dex plausible band, so they are stress tests, but the band itself is the calibration question. \emph{CRISTAL}: on the six dust-detected discs at the outer radius the flat law's median $\delta$ is $+0.106$ $[-0.138,+0.343]$ and the rival's $-0.203$ $[-0.477,+0.029]$; the route decides it, and two discs (CRISTAL-11 and -19, 0.35 and 0.47\,dex below the fit's baryon mass) decide the route. \emph{CFG228}: two of seven have $D<1$, a baryon-inventory inconsistency for every law, restored by the gas band; the pooled implied $s^*=4.33$ sits nearest the proxy's 4.6 and that coincidence carries no information. \emph{CFG229}: the seven sit at $g_{\rm bar}/a_0$ between 4.8 (ALESS 122.1) and 157, where flat and $H(z)$ differ by 0.005 to 0.09\,dex per galaxy; even with perfect baryons the pooled signal is below 2 SD, because the intrinsic scatter, velocity and inclination errors set the floor; ALESS 122.1 is one dispersion-supported galaxy whose envelope over every band corner and variant is 3.7 to 17.9 in $s^*$.
% AUDIT: B26

\textbf{Forecast lanes (no law scored).} CFG218: no sample is discriminating; the baryon-mass calibration needed for a 3$\sigma$-equivalent separation is 0.02 to 0.11\,dex (MUSE-DARK is systematic-limited: signal 0.092 against a systematic of 0.300). CFG218 read CRISTAL as limited by statistics (about 13 discs would remove that limit); the standing page records that CFG219 supersedes this: ``the limit is the gas calibration, not N''. CFG219: with a shared gas-calibration systematic of 0.25\,dex (about 0.15 on baryons) flat against $H(z)$ never reaches 3$\sigma$ at $R_e$ for $N\le20$; the shared gas calibration must be known to about 0.2\,dex (about 0.1--0.15\,dex on the baryon mass) however many discs there are. CFG221: the frozen pooled $z>3.5$ rule is one-sided in practice, and a fair two-sided test needs the gas calibration to about 0.15\,dex with about 50 outer-radius discs. CFG222: on RC100 the proxy sits between the flat law and $H(z)$ in every quantity (slope $-0.067$ against $-0.029$ and $-0.092$; the baryon tilt that would make each consistent is $-0.18$, $-0.07$ and $-0.25$\,dex), so the proxy too is calibration-bound; it is not a test of $\Lambda$CDM.
% AUDIT: B27

\textbf{Outside this note's lane list.} The standing page records the earlier KURVS ($z\approx1.5$) lean toward the rival as ``fragile and calibration-bound'', and for MUSE-DARK the defensible wording is that ``MUSE-DARK III's a$_0$ rise is not recovered on SED M$_\star$ + H$_2$; the data on disk cannot say whether the fitted masses or the SED route are biased.'' Neither changes the picture here.
% AUDIT: B28

\section{The gas calibration}\label{sec:gas}

\textbf{What tracer agreement measures (CFG224, CFG224b).} At the conversions each source states, CO, [CI] and dust gas masses agree with a pooled-scale error $K=0.038$ at $z<0.6$ (Stripe82, 78 galaxies) and $K=0.128$ at $z\approx1.2$ (14 galaxies); above $z\approx1.6$ $K$ is not estimated, because only three galaxies carry stated multi-tracer masses and their per-galaxy disagreements reach 0.6\,dex. With the per-galaxy optimised conversions of Dunne et al.\ [D22] the tracer masses agree by construction, and the informative content is how the conversion factors vary with $z$. ACE's dust-to-CO ratio at $z\approx2.2$ [ACE] lies 0.67\,dex below Stripe82's relation at the frozen setting (0.21 to 0.67 depending on the metallicity slope).
% AUDIT: B29

\textbf{What the referee found (CFG238).} The independent referee reproduces 145 of 146 scored numbers, but the reading of them does not survive. (a)~$K$ is the error of a pooled mean scale, not a per-galaxy gas-mass error: the catalogue error is 0.083 to 0.111\,dex per galaxy. (b)~The fixed-luminosity comparison is an extrapolation: one of eight cases has a single overlapping $L_{\rm IR}$ bin; with sample fixed effects 3 of 9 slopes reach about 3$\sigma$; the test has 21--52\% power for a real 0.05\,dex change, so ``no drift larger than about 0.08 to 0.12 dex across z = 0 to 2.5'' is what the data say at 80\% power, not ``no drift''. (c)~ACE's 0.2--0.7\,dex is a bracket of prescriptions on an extrapolation in metallicity (span 1.09\,dex across 30 cells), not a measured $z$-evolution. (d)~A common-mode bias shared by all tracers is invisible to every diagnostic (a $+0.30$\,dex shift on all tracers leaves every pair difference unchanged) and cannot be bounded from anything on disk that is independent of the dynamics under test. An earlier, more optimistic reading of CFG224 was withdrawn after this referee and is not used here.
% AUDIT: B30

\textbf{The standing wording, quoted exactly.} ``multi-tracer gas is the right route, but its calibration at z $\ge$ 1.6 is not shown to reach the ${\sim}$0.1 dex a decisive a$_0$(z) test needs; quote the gas band as a prescription bracket (${\sim}$0.2--0.7 dex at z $\approx$ 2.2) unless a common-mode anchor exists.''
% AUDIT: B31

\section{No dynamics-independent anchor}\label{sec:anchor}

A common-mode calibration needs an anchor that does not use the dynamics under test. A scoping of the public routes (paper text only, no downloads; \fpath{data_assembly/GAS_ANCHOR_SCOPING_2026-09-30.md}) reaches this bottom line: ``No public route pins the common-mode gas calibration at z $\ge$ 1.6 to 0.10 dex.''
% AUDIT: B32

\textbf{The best dynamics-free route, checked.} Heintz \& Watson [HW20] derive $\alpha_{\rm [CI]}(Z)$ from absorption-line column ratios in 19 GRB and QSO absorbers at $z$ 1.962--3.376, with a stated scatter of 0.2\,dex about $\log\alpha_{\rm [CI]}=-1.13\log(Z/Z_\odot)+1.33$. In our transcription of their Table~1 (made from the paper's PDF text, which is not committed to the repository): the 19 tabulated points scatter 0.64\,dex rms about that relation; 7 of the 19 $\alpha$ rows do not follow from the paper's own Eq.~3 applied to its columns; the $\alpha$ values of four QSO rows match a cyclic shift of the values computed from their sources (three agree to 0.00\,dex, the fourth is off by exactly 1.00 in the leading digit; the scoping note calls this a hypothesis, strongly supported, and it has not been put to the authors), and for the three GRB rows the $\alpha$ values follow the source CI$^*$ while the printed $N({\rm CI}^*)$ column is unexplained or matches no source value. Undoing the shift lowers the rms (our arithmetic, not the authors') from 0.64 to 0.55\,dex, not to the stated 0.2. Which version the authors fitted cannot be told from the paper. How the authors define their 0.2\,dex scatter (a fit with intrinsic scatter, or the error on the relation, would differ from an rms), and whether an erratum exists, were not checked. This is a transcription check of a published table, not a refutation of the relation.
% AUDIT: B33

\textbf{The other routes.} The absolute scale of Dunne et al.\ [D22] is assumed, not measured (a Milky-Way-like $\delta_{\rm GDR}=135$, i.e.\ $\kappa_H=1884$, or $X_{\rm CI}=1.6\times10^{-5}$ taken from [HW20]); it constrains the relative scales of CO, [CI] and dust to 0.11--0.15\,dex and says nothing on the common mode. The metallicity scale itself carries 0.14--0.18\,dex of calibration systematics [C19, B23]. For one object near solar metallicity the error budget is 0.35--0.7\,dex, 3 to 7 times the 0.10\,dex target (the scoping note's estimate). Only three Dunne et al.\ objects with CO, [CI] and dust also have a verifiable independent metallicity. Dynamical-mass inequalities are not anchors: 68\% of Stripe82 galaxies already violate $M_*+M_{\rm CO}+M_{\rm HI}\le M_{\rm dyn}$ (an aperture mismatch), and at $z$ 2--4 the dynamical masses are the output of the kinematic model under test (CFG238).
% AUDIT: B34

\section{The differential route: a within-survey redshift split}\label{sec:diff}

A split by redshift inside one survey cancels any calibration that is common to both halves. It does not cancel a calibration that drifts with redshift, and in the deep regime a drift is indistinguishable from an $a_0$ drift.

\textbf{KiDS-1000 by lens redshift (CFG255).} On 181,477 KiDS-bright lenses (on-disk per-lens sums; nothing fetched) the extreme photo-$z$ thirds have median $z$ 0.20 against 0.40 (late class) and 0.23 against 0.38 (early). Over that lever the rival predicts an amplitude change of $+0.018$\,dex and the flat law $+0.001$; the jackknife error on the combined amplitude is 0.038\,dex; the power is $\Delta\chi^2_{\rm pred}=0.14$ (canonical) / 0.17 (alternative) against the 9 required, so the pre-flight said NOT POSSIBLE before the measurement. The measurement is consistent with every line: $\chi^2/14$ is 19.8 for zero and flat, 19.2 for the rival and 15.1 for the $\Lambda$CDM proxy (all $p>0.1$); the amplitude is $+0.060\pm0.038$\,dex, noise-level and not read as evidence either way. With the rival's predicted ratio injected (MUTATE), flat is still not rejected ($p=0.10$): NON-DISCRIMINATING. An independent re-run reproduced every committed results JSON byte for byte.
% AUDIT: B35

\textbf{Why it is walled even at infinite statistics.} In the deep regime a stellar-mass calibration drift $\delta$ between the bins moves the amplitude by $\delta/2$, as an $a_0$ drift does (to within 6.1\% in the K1 bins, exactly in the deep limit), so the rival's 0.018\,dex needs the KiDS photometric $M_*$ stable to better than 0.036\,dex between $z\approx0.2$ and 0.4; no source quantifies that stability. What would make the test possible: lenses reaching $z\approx0.8$--1.0, where the rival's shift is 0.06--0.09\,dex; a per-split amplitude error $\lesssim0.02$\,dex; and a stellar-mass drift demonstrated below about 0.05\,dex. The data scoping found no public isolated-lens sample with stellar masses beyond $z=0.5$.
% AUDIT: B36

\textbf{The prior lane agrees.} An earlier, independent two-bin KiDS lensing fit with its own interpolation function (\fpath{real_research/reviews/lensing_rar/A0Z_LENSING_ZBIN_2026.md}; KiDS-1000 lensing RAR of [B21]) found $a_0({\rm high})/a_0({\rm low})=1.309\pm0.210$ (statistical only) at $z_{\rm eff}=0.2357$ and 0.3723, against a predicted 1.083 for the rising $E(z)$ branch (with that note's own $\Omega_m$; $\Omega_m=0.3$ gives 1.080, CFG241) and 1.000 for a constant $a_0$: nothing is excluded at $\ge2\sigma$. It adds a systematic CFG255's budget did not name: the magnitude-limited halves differ by 0.465\,dex in mean $\log M_{\rm gal}$, so a mass-dependent baryon term does not cancel between the halves; CFG255's model stacks follow each subset's own lens masses, but the gas-fraction prescription itself is uncertain. That strengthens NOT POSSIBLE; it is a null, not a detection.
% AUDIT: B37

\section{What would decide it}\label{sec:decide}

\textbf{The requirement.} The lanes land on the same order, 0.02 to 0.15\,dex. The ALMA pre-flight (CFG228, post hoc arithmetic on its own mocks) shows that the six ALPINE rotators \emph{would} separate flat from $H(z)$ if the baryon errors and the band were both $\le0.10$\,dex (signal 0.277 against $S_{\max}=0.078$ plus 2\,SD $=0.168$, with SD $=0.084$), and would not at 0.20\,dex. The other lanes give the same order: 0.05 to 0.10\,dex (CFG223), 0.02 to 0.11\,dex (CFG218), about 0.15\,dex on the gas with about 50 discs (CFG221), and a differential scale known to about $\pm0.06$\,dex (CFG217). The standing page's general statement: the 0.25-dex mass-scale floor caps the $z\approx2.5$ test at 2.3$\sigma$ against $a_0\propto H(z)$ and 1.3$\sigma$ against $\Lambda$CDM-native, whatever the sample size, so the calibration must reach about 0.1\,dex.
% AUDIT: B38

\textbf{What else a sample needs.} From the ALMA and class-M lanes: per-galaxy optimised multi-tracer gas with stellar masses that carry a measured error; inclinations known to 5$^\circ$ or better (4$^\circ$ to 31$^\circ$ in the ALPINE six); rotation curves resolved beyond 2\,$R_e$ with $\sigma(R)$ measured (the pressure term decides the sign of every $\delta$); an intrinsic scatter below about 0.05\,dex; and radii at which $g_{\rm bar}/a_0$ is 1 to 3 (for the class-M discs $g_{\rm bar}/a_0\approx3$ is reached only at 11 to 20\,kpc, 1.3 to 8 times the stated radii). None of the public samples on disk meets these together. A common-mode gas anchor at $z\ge1.6$ that does not use the dynamics would be needed in addition (Section~\ref{sec:anchor}).
% AUDIT: B39

\textbf{The near-term test that can kill B (and the bare law) is elsewhere.} Gaia DR4 (2 Dec 2026) tests the programme at $z=0$ with wide binaries; it does not measure $a_0(z)$. Candidate B (the law with hierarchical ownership) predicts $\hat\gamma=1.000$ and the bare law's floor is 1.161 (alt 1.192); with the frozen $\sigma_{\rm sys}=0.02$, 3$\sigma$ needs about 4,300 pairs (alt 2,900) out of about 30,000 expected. It cannot separate B from $\Lambda$CDM or Newton, which both predict 1.000: for B it is a survival test, never a confirmation.
% AUDIT: B40

\section{Limitations and what is not claimed}\label{sec:lim}

\begin{itemize}
\item This note compiles committed lanes and introduces no calculation. The implied-$a_0$ lever is a post hoc diagnostic of CFG223; the CFG228 requirement ($\le0.10$\,dex) is post hoc arithmetic on that lane's pre-flight mocks; CFG219's calibration statements use a post hoc systematic-inclusive statistic, because its frozen one was blind to a shared offset.
\item All $a_0(z)$ rows rest on author decompositions (halo priors, prior-anchored baryon masses, the authors' pressure corrections) or on our thin-disc baryon models; an implied $a_0$ is the scale that makes a decomposition match the kernel, not a direct measurement. Several lanes are not blind to published kinematics (CFG222, CFG223, CFG227, CFG228, CFG229 disclose this). Referees re-derived RC100 (CFG233), CRISTAL (CFG234), the compilation (CFG237) and the gas calibration (CFG238); the corrections they forced are folded in above. RC100 numbers are from the committed input table; a re-run on the published Table~3 values (17 cells differ) moves the RC100 slopes quoted here by at most 0.003 and no verdict (CFG216 input correction). Four NOEMA3D CO masses sit 0.24 to 0.26\,dex, and a fifth 0.149\,dex, below the paper's stated recipe (a data-assembly finding not checked against the paper); raising them changes no label (CFG213 note).
\item This note compiles the programme's own committed lanes; it does not review the high-redshift RAR literature or $\Lambda$CDM simulation predictions (CFG222 notes that a fair $\Lambda$CDM test uses simulated galaxies).
\item The $\Lambda$CDM proxy is not $\Lambda$CDM. Nothing here says the data favour the framework, the rival or $\Lambda$CDM; a lean is not a detection; ``calibration-limited'', ``NOT POSSIBLE'', ``non-diagnostic'' and ``descriptive'' carry their lanes' meanings. $\kappa=\tfrac12$ is fitted; the theory is not claimed closed; the cold mass is still required and no dark-matter particle is introduced.
\item The scoping of anchors read paper text, not data; other routes may exist that it did not find. Its caveats on the [HW20] re-reading are stated with it in Section~\ref{sec:anchor}.
\end{itemize}
% AUDIT: B41

\section{Reproducibility}\label{sec:rep}

Every number carried by an audit row of \fpath{PAPER38_audit.py} is re-read from a committed file: the script checks each value against its committed source at git HEAD and checks that the value is printed in the tex block that carries its audit tag (anywhere in the block, not in a particular sentence), reports how many rows match, and has two mutation modes (one altered expected value; one altered tex literal) that must fail. Numbers quoted only in running text are not checked by the script (23\% of the numeric tokens of version 1.1 by the referee's count); they were checked by the referee passes of CFG241, not by the script. The journal volumes and pages of the references are not audited by the script. The sources the script reads are \fpath{campaign_fresh_gravity/STANDING_2026-09-29.md}, \fpath{STANDING_2026-09-28.md} and \fpath{LEDGER.md} in the same directory, the READMEs and outputs of lanes CFG213 to CFG224 (with the CFG216 input-correction note), CFG227 to CFG229, CFG238, CFG240, CFG241 and CFG255 under \fpath{campaign_fresh_gravity/}, \fpath{data_assembly/GAS_ANCHOR_SCOPING_2026-09-30.md} and \fpath{real_research/reviews/lensing_rar/A0Z_LENSING_ZBIN_2026.md}; the referee READMEs of CFG233, CFG234 and CFG237 cited in Section~\ref{sec:lim} are not read by the script.

\textbf{Changes in version 1.2.} Wording corrections after the independent referee CFG241 (23 findings: 0 CRITICAL, 2 MAJOR, 13 MINOR, 8 NIT); no result changes and no source value changes. F-01: the abstract's ``binding limit in every lane'' is restated as the limit that survives an unlimited sample, and the lanes limited at present by statistics or the Newtonian regime are named. F-02: the [HW20] re-reading is stated as our transcription, the cyclic shift as the scoping note's hypothesis not put to the authors, the rms values as our arithmetic, with the unchecked scatter definition and erratum, in the abstract and in Section~\ref{sec:anchor} (the former Section~\ref{sec:lim} caveat moved there). F-03: the step dependence of the baryon lever, with the $\pm0.01$\,dex range. F-04: per-point medians against the per-galaxy $y$ spread, and the formal leverage of the real spread. F-05: the premises of the T4 floor. F-06: the $y$ range behind 0.117. F-07: the $\delta/2$ wall separated from the statistical verdicts. F-08: ``the same order'' for ``the same number''. F-09: DR4 is the near-term test B can only survive or fail, not ``decisive''. F-10: $E(z)$, $\Omega_m$ and the kernel stated. F-11: what the audit covers. F-12: the RC100 input correction and the NOEMA3D CO note. F-13: B23 and C19 journal data and the reference-check statement. F-14: a mislabelled threshold (the source gives SD $=0.084$, so 2\,SD $=0.168$; the conclusion is unchanged). F-15: the literature scope and the rival's attribution. N-01 to N-08: the flat law shared with standard MOND; the CFG219 quotation attributed to the standing page; the $\delta/2$ equality qualified; the Newtonian end of the determinant bound; audit tags renumbered B01 to B43 in text order (version 1.1 had no B34) and a hard-coded section number replaced; the scoping note's estimate labelled; CFG255's lens-mass clause restored; the full source list.
% AUDIT: B42

\section*{References}
\small
% v1.2 reference check (2026-10-01). Titles, first authors and arXiv identifiers of all nine arXiv entries were checked against the arXiv abstract records
% (export.arxiv.org API). Journal data: B21 against the arXiv journal_ref field and citations/REFERENCES.bib; DM14 and NS23 against citations/REFERENCES.bib only;
% C19 (MNRAS 485, 2092), D22 (MNRAS 517, 962) and HW20 (ApJL 889, L7) against the Crossref records of the DOIs that their arXiv records carry (C19's DOI resolver
% also redirects to mnras/article/485/2/2092); B23 (ApJ 958, 64) against the Crossref record found by exact title and first author, since its arXiv record carries no DOI;
% the ACE pair are arXiv preprints ("submitted to A&A" per their arXiv comments). citations/REFERENCES.bib has entries only for B21, DM14 and NS23; it was not edited.
[ACE] M.~Solimano et al.\ 2026, ``ALMA Chemical Evolution (ACE) survey: the dust content of subsolar metallicity galaxies at cosmic noon'', submitted to A\&A, arXiv:2609.21040; N.~N.~Geesink et al.\ 2026, ``ALMA Chemical Evolution (ACE) survey: dust-to-gas ratios in sub-solar metallicity galaxies at cosmic noon'', submitted to A\&A, arXiv:2609.20926.\par
[B21] M.~M.~Brouwer et al.\ 2021, ``The weak lensing radial acceleration relation: constraining modified gravity and cold dark matter theories with KiDS-1000'', A\&A 650, A113, arXiv:2106.11677 (Brouwer et al.\ 2021).\par
[B23] J.~E.~Birkin et al.\ 2023, ``JWST's TEMPLATES for star formation: the first resolved gas-phase metallicity maps of dust-obscured star-forming galaxies at $z\sim4$'', ApJ 958, 64, doi:10.3847/1538-4357/acf712, arXiv:2307.10412.\par
[C19] R.~T.~Coogan et al.\ 2019, ``Suppressed CO emission and high G/D ratios in $z=2$ galaxies with sub-solar gas-phase metallicity'', MNRAS 485, 2092, doi:10.1093/mnras/stz409, arXiv:1901.06390.\par
[D22] L.~Dunne et al.\ 2022, ``Dust, CO and [CI]: cross-calibration of molecular gas mass tracers in metal-rich galaxies across cosmic time'', MNRAS 517, 962, doi:10.1093/mnras/stac2098, arXiv:2208.01622.\par
[DM14] A.~A.~Dutton \& A.~V.~Macci\`o 2014, ``Cold dark matter haloes in the Planck era: evolution of structural parameters for Einasto and NFW profiles'', MNRAS 441, 3359, arXiv:1402.7073 (Dutton \& Macci\`o 2014).\par
[HW20] K.~E.~Heintz \& D.~Watson 2020, ``Direct measurement of the [CI] luminosity to molecular gas mass conversion factor in high-redshift star-forming galaxies'', ApJL 889, L7, arXiv:2001.05770.\par
[NS23] A.~Nestor Shachar et al.\ 2023, ``RC100: rotation curves of 100 massive star-forming galaxies at $z=0.6$--2.5 reveal little dark matter on galactic scales'', ApJ 944, 78, arXiv:2209.12199 (Nestor Shachar et al.\ 2023).\par
% AUDIT: B43

\end{document}
